\section{One physical Gaussian through changes of arithmetic type}
\label{sec:v58-common-raw-law}
We apply flat pressure contact to the same complete raw return law.
The finite phase and damping sum remains necessary, but the moving
means and complex Hessians can now be removed from its Gaussian
factor uniformly in the radius. The resulting formula has the
physical return covariance even when the arithmetic type changes
with the observation count.

Keep all peak data and partition functions of
\eqref{eq:v41-peak-data}--\eqref{eq:v41-transition-kernel}. For the
real target $y=(k_1,k_2,u,n)$ define
\begin{align}
 \Phi_{m,R}(y)&=\frac1c\operatorname{Re}
          \sum_j e^{-ia_j(R)\cdot y}z_j(R)^mB_j(R),
                                      \label{eq:v58-damped-residue}\\
 \mathcal L^\circ_{m,R}(y)&=
         \Phi_{m,R}(y)g_{\Omega_R}
                       \left(\frac{y-m\mu_R}{\sqrt m}\right),
                                      \label{eq:v58-common-kernel}\\
 \alpha_{m,R}(k,n,u)&=c^{-1}\Phi_{m,R}(k_1,k_2,u,n).
                                      \label{eq:v58-return-residue}
\end{align}
The factor $c=91/(10000\pi)$ is fixed. Each scalar sum is real and
uniformly bounded, but it need not be nonnegative at finite count.
Every retained branch occurs with its original amplitude, damping,
phase and partition. The roof phase $e^{-i(a_j)_3u}$ is not discarded.
Different real lifts of torus frequencies give the same expression at
the exact integer labels.

\subsection{Uniform removal of moving Gaussian data}
\begin{theorem}[Common Gaussian shape of the arithmetic transition]
\label{thm:v58-common-transition}
There are $a>0$ and a deterministic $\theta_m\to0$ such that
\begin{equation}\label{eq:v58-weighted-kernel}
 \sup_{R,y}e^{a|y-m\mu_R|^2/m}
      |\mathcal L_{m,R}(y)-\mathcal L^\circ_{m,R}(y)|
                                         \le\theta_m.
\end{equation}
The supremum may be taken over every real $y\in\R^4$ for a fixed
choice of frequency lifts; in the physical law its first, second and
fourth coordinates have their exact integer values. One admissible
bound, for all sufficiently large $m$, is
\begin{equation}\label{eq:v58-lorentz-modulus}
 \theta_m=C\{\omega_*(m^{-1/2})+
        \sqrt{\omega_*(m^{-1/2})}+m^{-1/6}+e^{-\sqrt m/4}\},
\end{equation}
where
$\omega_*(\delta)=\sup_{j,R:\,\kappa_j\le\delta}\|H_j-\Omega_R\|$.
The modulus is uniform over the whole radius family. In particular
no polynomial rate in $m$ is asserted.
\end{theorem}
\begin{proof}
In Corollary~\ref{cor:v58-common-shape} take $d=4$,
$Z=(y-m\mu_R)/\sqrt m$, $v_j=\mu_j-\mu_R$ and
\[
 q_{j,m}(R,y)=c^{-1}e^{-ia_j(R)\cdot y}z_j(R)^mB_j(R).
\]
The $B_j$ are uniformly bounded and $|z_j|^m=e^{-m\kappa_j}$.
Corollary~\ref{cor:v58-lorentz-contact} verifies the hypotheses.
Taking real parts gives \eqref{eq:v58-weighted-kernel}, and
$\delta=m^{-1/2}$ gives \eqref{eq:v58-lorentz-modulus}.
This uses only the fixed peak charts. It substitutes no growing
roof band into a spectral estimate.
\end{proof}

\begin{proposition}[Slow roof variation and asymptotic reference positivity]
\label{prop:v58-reference-positivity}
For the same exact integer labels $1\le n\le m$, $k\in\Z^2$,
\begin{equation}\label{eq:v58-reference-negative}
 \sup_{R,n,k,u}(\mathcal L^\circ_{m,R}(k_1,k_2,u,n))_-\to0.
\end{equation}
Consequently, for every fixed $M<\infty$,
\[
 \sup_{R,n,k,u:\,|y-m\mu_R|\le M\sqrt m}
                  (\Phi_{m,R}(y))_-\to0.
\]
These assertions concern the reference, not an upper bound for the
positive original density. The finite reference is still treated as
signed.
\end{proposition}
\begin{proof}
Put $b_j=(a_j)_3$. At every zero of $\kappa_j$ one has $b_j=0$
by Lemma~\ref{lem:v41-moving-maxima}. Continuity and compactness
therefore give
\[
 \beta(\delta):=\sup_{j,R:\,\kappa_j\le\delta}|b_j|\to0.
\]
Differentiating the finite sum, and using uniform Gaussian derivative
bounds, shows
\begin{equation}\label{eq:v58-roof-lipschitz}
 \sup_{R,y}|\partial_u\mathcal L^\circ_{m,R}(y)|
 \le C\left(m^{-1/2}+\sum_j|b_j|e^{-m\kappa_j}\right)=:d_m\to0.
\end{equation}
For example the sum is bounded by $C\{\beta(m^{-1/2})+
 e^{-\sqrt m}\}$. The coefficients outside the Gaussian do not
have to be independent of the roof for this argument.

Let $e_m\to0$ bound the inherited complete local variation error on
every unit roof interval from Theorem~\ref{thm:v49-raw-local-TV}.
For $P=m^2p\ge0$ and $G^\circ=\mathcal L^\circ$, the new kernel
comparison gives
\[
 \int_u^{u+1}|P(v)-G^\circ(v)|\,dv\le e_m+\theta_m.
\]
If $G^\circ(u)=-h<0$, then
$G^\circ(v)\le-h+d_m$ for $u\le v\le u+1$.
When $h>d_m$, positivity of $P$ bounds the last integral below by
$h-d_m$. Hence $(G^\circ(u))_-\le e_m+\theta_m+d_m$ at every
roof; the case $h\le d_m$ is immediate. The density is extended by
zero outside its physical roof support, as in the local theorem.
Finally the Gaussian $g_{\Omega_R}$ has a positive uniform lower
bound on $|Z|\le M$, proving the assertion for $\Phi$.
No estimate of the positive part of $P-G^\circ$ is inferred.
\end{proof}

\subsection{The fixed-return covariance on the entire record space}
Retain $\mathcal O_n$, $d\zeta$ and $V_{n,R}$ from
Section~\ref{sec:v57-transition-tails}. Define the signed density
\begin{equation}\label{eq:v58-physical-reference}
 q^\circ_{n,R}(k,m,u)=n^{-2}\alpha_{m,R}(k,n,u)
                                  g_{D_R}(V_{n,R}(k,m,u)),
                  \quad(k,m,u)\in\mathcal O_n.
\end{equation}
All collision counts $m\ge n$, displacements $k$ and positive roofs
are present. Define $Z^\circ_{n,R}=\int(q^\circ_{n,R})_+d\zeta$ and,
when this number is positive, let $\mathsf A^\circ_{n,R}$ be the
probability with density $(q^\circ_{n,R})_+/Z^\circ_{n,R}$.

\begin{theorem}[The uniform arithmetic Gaussian law of the full record]
\label{thm:v58-canonical-raw}
For the original actual return record, uniformly in the radius,
\begin{equation}\label{eq:v58-canonical-L1}
 E_n^\circ:=\sup_R\sum_{m=n}^\infty\sum_{k\in\Z^2}
     \int_0^\infty|p_{n,R}(k,m,u)-q^\circ_{n,R}(k,m,u)|\,du
                                      \longrightarrow0.
\end{equation}
Its reference satisfies $Z^\circ_{n,R}\to1$ uniformly and
\[
 d_{\rm TV}(\mathsf P_{n,R},\mathsf A^\circ_{n,R})\le E_n^\circ.
\]
The negative mass of $q^\circ$ tends to zero. Here $d_{\rm TV}$ for
probabilities is one half of variation mass.
At each fixed radius the factor $\alpha_{m,R}$ may be replaced by
$\mathfrak a_R(k,n,m)$, including its zero classes. Uniformly through
arithmetic changes the finite damped factor is retained.
\end{theorem}
\begin{proof}
First compare the original transition density
$\gamma_{n,R}=m^{-2}\mathcal L_{m,R}$ with
$m^{-2}\mathcal L^\circ_{m,R}$. Equation
\eqref{eq:v58-weighted-kernel} and the Gaussian count sum
\eqref{eq:v57-count-tail-sum} give directly
\begin{equation}\label{eq:v58-first-global-replacement}
 \int_{\mathcal O_n}m^{-2}|\mathcal L_{m,R}-\mathcal L^\circ_{m,R}|
                    \,d\zeta\le C\sup_{m\ge n}\theta_m\to0.
\end{equation}
Both kernels have uniform reference tails by the same estimate.

We next change from the collision scale to the return scale. On the
fixed set $|V_{n,R}|\le M$ one has $m=n/c+O_M(\sqrt n)$.
The linear relations in \eqref{eq:v37-covariance-change} give
\[
 \frac{y-m\mu_R}{\sqrt m}=\sqrt{n/m}\,L_RV_{n,R},\qquad
 \Omega_R=cL_RD_RL_R^{\mathsf T},\qquad \det L_R=c.
\]
Thus
$g_{\Omega_R}(\sqrt cL_RV)=c^{-3}g_{D_R}(V)$.
Together with $n^2/m^2\to c^2$, uniformly on this set, this proves
\begin{equation}\label{eq:v58-covariance-scale}
 \sup_{R,\,|V_{n,R}|\le M}
 n^2|m^{-2}\mathcal L^\circ_{m,R}(y)-q^\circ_{n,R}(k,m,u)|\to0.
\end{equation}
The remaining factor is $\Phi/c=\alpha$, precisely as in
\eqref{eq:v58-return-residue}; the section normalization is not
applied a second time.

Here is the count bookkeeping for integration. On $|V|\le M$ there
are $O_M(n)$ displacement labels, $O_M(\sqrt n)$ collision counts,
and $O_M(\sqrt n)$ unit roof intervals. Therefore the bound in
\eqref{eq:v58-covariance-scale} integrates with the factor
\begin{equation}\label{eq:v58-cell-bookkeeping}
 \underbrace{O_M(n)\,O_M(\sqrt n)\,O_M(\sqrt n)}_{O_M(n^2)
                       \ \mathrm{unit\ cells}}
       \underbrace{O_M(m^{-2})}_{O_M(n^{-2})}\,o(1)=o(1).
\end{equation}
Equivalently use its displayed $n^{-2}$ normalization. All constants
here are for a fixed $M$; no count index has been averaged.

Since $|\alpha|\le C$ and $D_R$ is uniformly elliptic, direct
Gaussian summation on the three scaled lattice coordinates and
integration on the roof give
\[
 \int_{\mathcal O_n}|q^\circ_{n,R}|\,d\zeta\le C,\qquad
 \int_{|V_{n,R}|>M}|q^\circ_{n,R}|\,d\zeta\le Ce^{-aM^2}.
\]
For example sum $e^{-a k_i^2/n}$ and
$e^{-a(m-n/c)^2/n}$ separately, each at cost $C\sqrt n$, and
integrate the roof Gaussian at the same cost. The prefactor is
$n^{-2}$. Spend half the exponent for the tail. The old reference
has the separate tail bound \eqref{eq:v57-reference-tail}.
Let $n\to\infty$ first and then $M\to\infty$. Equations
\eqref{eq:v58-first-global-replacement}--\eqref{eq:v58-covariance-scale}
prove $\sup_R\int|\gamma-q^\circ|d\zeta\to0$.
Theorem~\ref{thm:v57-global-raw-TV} now proves
\eqref{eq:v58-canonical-L1}.

For $p\ge0$ one has $|p-q_+^\circ|\le|p-q^\circ|$, so
$|Z^\circ-1|\le E_n^\circ$ and normalization costs at most a second
copy of this error in variation mass. Division by two proves the
probability bound. On $q^\circ<0$ its negative part is at most
$|p-q^\circ|$, proving the negative-mass assertion.

At fixed $R$, all peaks with positive damping decay exponentially
in $m$; the remaining peaks are actual resonances with zero roof
coordinate. Their original phase sum, divided by $c^2$, is exactly
$\mathfrak a_R(k,n,m)$, by the fixed-radius specialization in
Section~\ref{sec:v41-transition}. The difference of the two bounded
coefficients tends to zero uniformly in the labels and roof. The
same Gaussian summation proves the fixed-radius assertion on all of
$\mathcal O_n$. It does not replace a zero residue by a positive
conditional law.
\end{proof}

\subsection{The same coupled paths and exact observation event}
Let $\mathsf H_R$ and the observation-variation/path-BL norm have
exactly the definitions in Section~\ref{sec:v57-coupled-bridges}.
The pair consists of the collision bridge and the actual-return
bridge from the \emph{same} trajectory; its reference is the graph
law $\operatorname{Law}(\mathbb B_{\Omega_R},
 c^{-1/2}L_R^{-1}\mathbb B_{\Omega_R})$.

\begin{corollary}[Canonical reference for the coupled record and paths]
\label{cor:v58-canonical-bridge}
With $\mathsf A^\circ_{n,R}$ in place of $\mathsf A_{n,R}$,
\begin{equation}\label{eq:v58-joint-canonical}
 \Delta_n^\circ:=\sup_R
 \left\|\operatorname{Law}_{\nu_R^*}
 (J_{n,R},\mathcal B_{N_{n,R},R},\mathcal Y_{n,R})
           -\mathsf A^\circ_{n,R}\otimes\mathsf H_R
                  \right\|_{\mathrm{obs},\mathrm{BL}}\to0.
\end{equation}
The two bridges inside $\mathsf H_R$ are not independent. For any
record-only Markov observation kernel and any event in its complete
output with reference probability at least $r_n>\Delta_n^\circ$,
the conditional output probability-TV error is at most
$\Delta_n^\circ/(r_n-\Delta_n^\circ)$, and the joint path-dual error
is at most $2\Delta_n^\circ/(r_n-\Delta_n^\circ)$.
Thus convergence after selection requires
$\Delta_n^\circ/r_n\to0$. A path-reading selector is not covered.
\end{corollary}
\begin{proof}
The preceding theorem and Theorem~\ref{thm:v57-global-raw-TV} give
$d_{\rm TV}(\mathsf A_{n,R},\mathsf A^\circ_{n,R})\to0$ uniformly.
Tensoring a signed observation measure with the probability
$\mathsf H_R$ gives mixed norm at most its variation mass. Apply
Theorem~\ref{thm:v57-global-coupled} and the triangle inequality to
obtain \eqref{eq:v58-joint-canonical}. Pullback through a record-only
kernel preserves the path test ball, and restriction to the full
output event does not increase the unnormalized error. The difference
of event masses is at most $\Delta_n^\circ$; normalization costs at
most twice that error divided by the true mass, which is at least
$r_n-\Delta_n^\circ$. On the scalar marginal divide by two for
probability TV. This reproduces the stated constants without changing
the event or inferring a convergence rate for a microscopic selection.
\end{proof}

\paragraph{The original pointwise endpoint.}
The common Gaussian shape removes moving centers and complex
covariances from the reference, not positive incidence or clearance
from the source. The one-sided raw identity still has those two
positive source densities at the exact labels. Their essential heights
are not bounded by Theorem~\ref{thm:v58-canonical-raw} or by the
reference positivity proposition. In particular neither unrestricted
same-roof bridges nor forward essential likelihoods are asserted here.
The original pointwise criterion, arithmetic zero classes, half-open
occupation at $0,\ldots,m-1$ and clearance at collision $j+1$ remain
unchanged.
